\subsection{Lineare Gleichungen und lineare Gleichungssysteme}
\label{subsec:linear_systems_motivation_and_definition}

\begin{DefinitionBox}{Lineare Gleichung}{linear_systems_linear_equation}
  Eine Gleichung in den Unbekannten $x_1,\dots,x_n$ heißt \emph{linear}, wenn
  sie die Form
  \[
    a_1x_1+a_2x_2+\dots+a_nx_n=b
  \]
  besitzt. Die Zahlen $a_1,\dots,a_n$ heißen \emph{Koeffizienten}, $b$ heißt
  \emph{Absolutglied}.
\end{DefinitionBox}

Die Unbekannten werden in einer linearen Gleichung weder miteinander
multipliziert noch potenziert und stehen nicht in nichtlinearen Funktionen.
Konstanten dürfen dagegen beliebig dargestellt sein:

\begin{ExampleBox}{Linear oder nichtlinear?}{linear_systems_linearity_examples}
  \begin{enumerate}[label=\alph*),itemsep=0.6em]
    \item $2x_1+3x_2=2^2$ ist \Cloze{3.2cm}{linear}.
    \item $x_1x_2=1$ ist \Cloze{3.2cm}{nicht linear}.
    \item $x_1^2+\sin(x_2)=1$ ist \Cloze{3.2cm}{nicht linear}.
  \end{enumerate}
\end{ExampleBox}

\begin{DefinitionBox}{Lineares Gleichungssystem}{linear_systems_system}
  Ein \emph{lineares Gleichungssystem} (LGS) aus $m$ Gleichungen in
  $n$ Unbekannten besitzt die Form
  \[
    \begin{aligned}
      a_{11}x_1+a_{12}x_2+\dots+a_{1n}x_n&=b_1,\\
      a_{21}x_1+a_{22}x_2+\dots+a_{2n}x_n&=b_2,\\
      &\mathrel{\phantom{=}}\vdots\\
      a_{m1}x_1+a_{m2}x_2+\dots+a_{mn}x_n&=b_m.
    \end{aligned}
  \]
  Man spricht kurz von einem $m\times n$-LGS.
\end{DefinitionBox}
